\subsection{Penalized Linear Estimation}
Regularisation controls model complexity by penalising unstable parameter directions or by constraining the feasible set.

\begin{definition} \label{def:sl-penalized-linear-estimator} \hlt{Penalised Linear Estimation}\\
For empirical loss $L_n(\beta)$, a penalised estimator solves $\widehat{\beta}_{\lambda}\in \operatorname*{arg\,min}_{\beta} \{L_n(\beta)+\lambda\Omega(\beta)\}$.\\
The choices $\Omega(\beta)=\norm{\beta}_2^2$ and $\Omega(\beta)=\norm{\beta}_1$ give ridge and lasso penalties.
\end{definition}

\begin{definition} \label{def:sl-subset-ridge-lasso-linear-models} \hlt{Subset Selection, Ridge, and Lasso}\\
Let $y\in\R^n$ be centered and let $X\in\R^{n\times p}$ have centered, standardised columns.
\begin{enumerate}[label=(\roman*)]
\setlength{\itemsep}{0pt}
\item \hlt{Best subset}: for a size $m$, solve $\widehat{\beta}^{\mathrm{sub}}_m\in \operatorname*{arg\,min}\limits_{|\operatorname{supp}(\beta)|\leq m}\norm{y-X\beta}_2^2$
\item \hlt{Ridge}: for $\lambda\geq0$, solve $\widehat{\beta}^{\mathrm{ridge}}_\lambda\in \operatorname*{arg\,min}\limits_{\beta} \{\norm{y-X\beta}_2^2+\lambda\norm{\beta}_2^2\}$
\item \hlt{Lasso}: for $\lambda\geq0$, solve $\widehat{\beta}^{\mathrm{lasso}}_\lambda\in \operatorname*{arg\,min}\limits_{\beta} \left\{\frac12\norm{y-X\beta}_2^2+\lambda\norm{\beta}_1\right\}$
\end{enumerate}
Centering separates the unpenalised intercept, estimated by the mean on the raw scale.
\end{definition}

\begin{remark} \label{rem:sl-ridge-lasso-geometry} \hlt{Penalty Geometry}\\
Assume $X^TX=I$ and set $z_j=x_j^Ty$. For subset size $m$, best subset selection keeps the $m$ largest $|z_j|$.\\
Ridge and lasso have coordinates $\widehat{\beta}^{\mathrm{ridge}}_{\lambda,j}=\frac{z_j}{1+\lambda}$ and $\widehat{\beta}^{\mathrm{lasso}}_{\lambda,j}=\sign(z_j)(|z_j|-\lambda)_+$.\\
Thus $\ell^2$ shrinkage is smooth, while $\ell^1$ shrinkage can set coefficients exactly to zero.
\end{remark}

\begin{figure}[H]
\centering
% Schematic redraw for ridge/lasso constraint geometry.
\begin{tikzpicture}[
  x=1cm,
  y=1cm,
  >=Stealth,
  line cap=round,
  line join=round,
  every node/.style={font=\scriptsize},
  axis/.style={->, gray!70, line width=0.45pt},
  contour/.style={red!65!black, line width=0.45pt},
  constraint/.style={draw=cyan!75!black, fill=cyan!16, line width=0.75pt},
  solution/.style={circle, fill=blue!70!black, inner sep=0pt, minimum size=2.4pt},
  ls/.style={circle, fill=black, inner sep=0pt, minimum size=2.4pt}
]
  \path[use as bounding box] (-1.55,-1.18) rectangle (7.45,1.95);

  \begin{scope}[shift={(1.0,0)}]
    \draw[axis] (-1.35,0) -- (1.45,0) node[right] {$\beta_1$};
    \draw[axis] (0,-1.25) -- (0,1.45) node[above] {$\beta_2$};
    \filldraw[constraint] (0,1.0) -- (1.0,0) -- (0,-1.0) -- (-1.0,0) -- cycle;
    \foreach \rx/\ry in {1.35/0.45,0.98/0.32,0.63/0.20}
      \draw[contour, rotate around={31:(1.10,0.92)}] (1.10,0.92) ellipse [x radius=\rx, y radius=\ry];
    \node[ls] at (1.10,0.92) {};
    \node[solution] at (1.0,0) {};
  \end{scope}

  \begin{scope}[shift={(5.65,0)}]
    \draw[axis] (-1.35,0) -- (1.45,0) node[right] {$\beta_1$};
    \draw[axis] (0,-1.25) -- (0,1.45) node[above] {$\beta_2$};
    \filldraw[constraint] (0,0) circle (1.0);
    \foreach \rx/\ry in {1.35/0.45,0.98/0.32,0.63/0.20}
      \draw[contour, rotate around={31:(1.10,0.92)}] (1.10,0.92) ellipse [x radius=\rx, y radius=\ry];
    \node[ls] at (1.10,0.92) {};
    \node[solution] at (0.74,0.67) {};
  \end{scope}
\end{tikzpicture}

\caption{Left: lasso $\ell^1$ diamond; right: ridge $\ell^2$ disk.}
\label{fig:sl-ridge-lasso-constraint-geometry}
\end{figure}

\begin{proposition} \label{prop:sl-ridge-effective-degrees} \hlt{Ridge Shrinkage and Effective Degrees}\\
Let $X\in\R^{n\times p}$ be centered with rank $r$ and thin SVD $X=UDV^T$, where $D=\operatorname{diag}(d_1,\ldots,d_r)$ and $d_j>0$.\\
For $\lambda>0$, set $S_\lambda\defeq X(X^TX+\lambda I_p)^{-1}X^T$ and $s_j(\lambda)\defeq \frac{d_j^2}{d_j^2+\lambda}$. Then $S_\lambda=U\operatorname{diag}(s_1(\lambda),\ldots,s_r(\lambda))U^T$.\\
So $\widehat{y}^{\mathrm{ridge}}_\lambda\defeq X\widehat{\beta}^{\mathrm{ridge}}_\lambda=S_\lambda y=\sum\limits_{j=1}^r s_j(\lambda)u_ju_j^Ty$ and $\operatorname{df}(\lambda)\defeq\operatorname{tr}(S_\lambda)=\sum\limits_{j=1}^r s_j(\lambda)$.
\begin{proof} Since $X^TX=VD^2V^T$ and $X^T=VDU^T$, we have $V^T(X^TX+\lambda I_p)^{-1}V=(D^2+\lambda I_r)^{-1}$.\\
Hence $S_\lambda=UDV^T(X^TX+\lambda I_p)^{-1}VDU^T=UD(D^2+\lambda I_r)^{-1}DU^T=U\operatorname{diag}(s_1(\lambda),\ldots,s_r(\lambda))U^T$.\\
Apply to $y$ and take traces, using $U^TU=I_r$.\\
\hlt{Structural Commentary}: Ridge is diagonal in the left singular-vector basis of $X$.\\
\hlt{Key Insight}: Ridge keeps direction $u_j$ by $s_j(\lambda)$; weak directions count less in the fit and $\operatorname{df}(\lambda)$.
\end{proof}
\end{proposition}
